\documentclass[12pt, a4paper]{article}

\usepackage{amsmath, amssymb, amsfonts}
\usepackage{float}
\usepackage{graphicx}
\usepackage{listings}
\usepackage{xcolor}
\usepackage[linesnumbered, ruled]{algorithm2e}
\usepackage{subcaption}
\usepackage[hidelinks]{hyperref}
% Each algorithm is written twice: the algorithm2e block (what the PDF typesets)
% and a plain-text Verbatim block right after it (what the web shows — pandoc
% renders Verbatim as a copyable <pre> code block). \excludecomment makes pdflatex
% SKIP the Verbatim blocks, so the PDF shows only the algorithm2e version; pandoc
% conversely drops each algorithm2e to an (empty, CSS-hidden) div and keeps the
% Verbatim. Because pdflatex never sees the Verbatim, its Unicode maths glyphs
% (⟨⟩ ← τ ∈ ∞ …) need no inputenc/literate setup.
\usepackage{comment}
\excludecomment{Verbatim}

\definecolor{codegreen}{rgb}{0,0.6,0}
\definecolor{codegray}{rgb}{0.5,0.5,0.5}
\definecolor{codepurple}{rgb}{0.58,0,0.82}
\hypersetup{linktoc=all}

\lstset{
    language=Python,
    basicstyle=\ttfamily\small,
    keywordstyle=\color{blue},
    commentstyle=\color{codegreen},
    stringstyle=\color{codepurple},
    breaklines=true,
    numbers=left,
    numberstyle=\tiny\color{codegray},
    frame=single,
    showstringspaces=false
}

\title{AI Masterclass Test 1}
\author{}
\date{}


\begin{document}

\section{Game Trees - Backward Induction}

\begin{algorithm}[H]
\caption{Zermelo's Algorithm}
\DontPrintSemicolon

$done \leftarrow leaves((S, E))$\;
\While{$done \neq S$}{
    $next \leftarrow \{s' \in S \setminus done \mid children(s', (S, E)) \subseteq done\}$\;
    \For{$s \in next$}{
        $i \leftarrow owner(s)$\;
        $C \leftarrow children(s, (S, E))$\;
        $O \leftarrow \arg \max_{s' \in C} u_i(s')$ \tcp*[r]{optimal choices}
        $s'' \leftarrow$ any element of $O$\;
        \For{$j \in N$}{
            $u_j(s) \leftarrow u_j(s'')$ \tcp*[r]{back utilities up}
        }
        $done \leftarrow done \cup \{s\}$ \tcp*[r]{we have processed $s$}
    }
}
\end{algorithm}
\begin{Verbatim}
# Zermelo's Algorithm
done ← leaves((S, E))
while done ≠ S:
    next ← { s' ∈ S ∖ done | children(s', (S, E)) ⊆ done }
    for s ∈ next:
        i ← owner(s)
        C ← children(s, (S, E))
        O ← arg max_{s' ∈ C} uᵢ(s')          # optimal choices
        s'' ← any element of O
        for j ∈ N:
            uⱼ(s) ← uⱼ(s'')                   # back utilities up
        done ← done ∪ {s}                     # we have processed s
\end{Verbatim}

\medskip
\noindent Extensive form games are games that are played sequentially. It can be defined by
\[G = (N, (S, E, s_0), owner, a, u_1,\ldots,u_n)\]
\begin{itemize}
        \item $N = \{1,\ldots,n\}$ is the set of players
        \item $(S, E, s_0)$ is a finite tree with vertex set $S$,
        edge set $E \subset S \times S$, and root $s_0 \in S$.
        \item $owner: interior((S, E))\rightarrow N$ specifies the owner of each decision node
        \item $a: E \rightarrow A$ associates each edge $(s, s') \in E$ with an action
        \item $u_i : leaves((S, E))\rightarrow\mathbb{R}$ is $i$'s utility function.
\end{itemize}

\noindent Zermelo's Algorithm / backward induction can solve extensive form games
\begin{itemize}
    \item Is guaranteed to terminate
    \item runs in time polynomial in the size of the game tree, $O(b^d)$
    \item finds optimal strategies for both players
\end{itemize}


\noindent\textbf{Example} Centipede Game
\begin{figure}[H]
    \centering
    \includegraphics[width=0.4\textwidth]{images/centipede-game.png} 
    % \caption{$h(s)$ = number of misplaced tiles or sum of Manatten distances of each tile to their goal position}
    \label{fig:my_handwritten_notes}
\end{figure}
\noindent Based on Zermelo's Algorithm,
\begin{enumerate}
    \item At $v_3$, Player 2 chooses $R$, since 4>3.
    $v_3$ effectively becomes the payoff (1, 4).
    \item At $v_2$, Player 1 chooses $R$, since 2>1.
    $v_2$ effectively becomes the payoff (2, 2).
    \item At $v_1$, Player 2 chooses $R$, since 3>2.
    $v_1$ effectively becomes the payoff (0, 3).
    \item At $v_0$, Plyaer 1 chooses $R$, since 1>0.
    $v_0$ effectively becomes the payoff (1, 1).
\end{enumerate}
Based on backward induction, when both players are perfectly rational and can predict each other's moves,
Player 1 just chooses $R$ and stops the game.\newline
The paradox is that they could each have had a greater payoff of (3, 3) if they trusted each other (deviating from the rational path) 
and played $D$ all the way.

\end{document}
